\section{Conditional Support Expressions and Assigned Inputs}\label{sec:fusion-contact-ledger}

\subsection{Inherited fusion domain and the present calculation}\label{subsec:fusion-closure-derivation}
\paragraph{Setup.}
For admitted forms \(F,G\) with certificate-bound cycle-valued compressions \(\widehat F,\widehat G\), Main defines
\begin{equation}
\operatorname{Fuse}_B(F,G)
=\operatorname{Close}_B(\widehat F\,\mu\,\widehat G)
\label{eq:matter1-fusion-closure}
\end{equation}
and
\begin{equation}
RD_{\mathrm{fuse}}(F,G;B)
=\Pi^-_{\widehat F}(B)+H_{\mu,\widehat F,\widehat G}(B)+\Pi^+_{\widehat G}(B).
\label{eq:matter1-rd-fuse}
\end{equation}
Their exact predecessors are \mainequation{eq:cycle-valued-field-compression}, \mainequation{eq:fusion-closure} and \mainequation{eq:fusion-branch-hinge-duration}. These operators retain their domains. The present elemental table supplies no inner or outer fusion certificate and does not assert an instance of \eqnref{eq:matter1-fusion-closure}.

\paragraph{Derivation.}
The declared reference operands use the exact polynomial values in \tabref{tab:matter1-curling-curl-primitives}:
\begin{equation}
u_-=30Z+68(A-Z),\qquad h=1,\qquad u_+=4s.
\label{eq:matter1-rung-operands}
\end{equation}
Define their conditional support sum
\begin{equation}
U(Z,A,s)=u_-+h+u_+=30Z+68(A-Z)+1+4s.
\label{eq:matter1-support-sum}
\end{equation}
This preserves the scalar expression formerly stored as \texttt{rd\_fuse}. The active field \texttt{support\_sum} names its established arithmetic scope. Equality to a native reflection-duration accessor additionally requires the inherited instance prerequisites.

\paragraph{Specialization.}
The nuclear count pair \((Z,A-Z)\) in \eqnref{eq:matter1-rung-operands} is distinct from the local tuple \((a,b,c)\) in \eqnref{eq:matter1-seat-abc}. At fixed \(Z\), changing the supplied integer \(A\) leaves \(m,s,a,b,c\) unchanged and gives
\[
U(Z,A+\Delta A,s)-U(Z,A,s)=68\Delta A.
\]
For carbon, \(Z=s=6\), the conditional choices \(A=12,13,14\) give \(U=613,681,749\). This is an algebraic illustration, not a sourced isotope catalog, stability statement or decay/fusion transition.

The reference A values are carried from the reviewed candidate's input list without alteration. Their original per-element selection policy was not recorded. The input provenance identifies that candidate source and this limitation; no broad AME or NIST citation establishes each selected A. Reference A is an integer nucleon-count input, not standard atomic weight. Projected A values for slots 119--126 remain explicitly illustrative fixture operands, not observed nuclide records. The inputs are recorded in \path{manual2/data/m1/reference_inputs.csv}, \path{model_inputs.json} and \path{input_provenance.json} in the same directory.

\paragraph{Falsification.}
A change in A that changes the grouping or local tuple violates this specialization. Calling a numerically correct \(U\) a certified fusion duration without its input and outer certificates exceeds its domain.

\aodprovenance{\mainequation{eq:curling-support-model}, \mainequation{eq:fusion-branch-hinge-duration}, \eqnref{eq:matter1-rung-operands}, \eqnref{eq:matter1-seat-abc}.}

\subsection{Assigned classifications, residuals and occupancy}\label{subsec:matter1-assigned-inputs}
\paragraph{Setup.}
The fixture keeps the candidate's assignment sets:
\[
\mathcal A=\{1,8,9,16,17,34,35,52,53\},\qquad
\mathcal X=\{43,61\}\cup\{84,\ldots,118\}.
\]
They define assigned acceptor and assigned shedding classes. They are input classifications; they are not derived from \(Z,A\), parity, or an executed native calculation. Let \(L=2m\), \(C^{\mathrm{fix}}_{\mathrm{close}}=3L\), and \(w^{\mathrm{fix}}=L\).

\paragraph{Derivation.}
With assigned pressure \(P^{\mathrm{fix}}\), retain the residual arithmetic
\begin{equation}
S^{\mathrm{fix}}=P^{\mathrm{fix}}-C^{\mathrm{fix}}_{\mathrm{close}},
\qquad
X^{\mathrm{fix}}=\max(0,S^{\mathrm{fix}}).
\label{eq:matter1-rung-residual-shedding}
\end{equation}
The pressure input is \(C^{\mathrm{fix}}_{\mathrm{close}}-3\) on \(\mathcal A\), equal to \(C^{\mathrm{fix}}_{\mathrm{close}}\) on the remaining labelled neutral rows, and \(C^{\mathrm{fix}}_{\mathrm{close}}+3u_Z\) on \(\mathcal X\), where
\[
u_Z=1\ \ (Z=43,61),\qquad
u_Z=\lfloor(Z-84)/6\rfloor+1\ \ (84\le Z\le118).
\]
For illustrative unfilled slots it is \(C^{\mathrm{fix}}_{\mathrm{close}}+3(Z-118)\). Assigned duration values are respectively \(L\), \(L+2\), \(L-1\), and \(0\). Assigned route and Yro labels are carried as fixture data. These values do not demonstrate retention, instability, transfer or non-retention.

The arithmetic uses the form of \mainequation{eq:closure-residual} and \mainequation{eq:shedding-excess}; it does not supply their native pressure, field and window evidence. Likewise, a declared route label does not execute the sheddic partition, lift or higher-scope fusion in \mainequation{eq:shedding-route-partition}, \mainequation{eq:exoshedding-lift} and \mainequation{eq:higher-scope-fusion}.

\paragraph{Specialization.}
The 118 labelled rows have disjoint assigned residual classes: 72 neutral, 9 deficit/acceptor and 37 positive-surplus/shedding. Boundary position is a separate attribute: eight labelled endpoints are \(\{6,14,24,36,50,66,84,104\}\). Po and Rf are both endpoints and assigned shedding rows. Adding endpoint counts to route-class counts would double-count those attributes. Of the 72 neutral rows, six are endpoints and 66 are not.

The \(\mathsf C_{22}\) fraction \(7/11\) records 14 labelled slots out of 22. Slots 119--126 are unfilled in this reference; their assigned fixture values establish neither a physical cutoff nor an impossibility of further elements.

\paragraph{Falsification.}
The assignment table and its equations must agree exactly. Positive \(X^{\mathrm{fix}}\) for negative \(S^{\mathrm{fix}}\), or a changed assigned set without a new input binding, fails the fixture. Agreement with these definitions is not empirical validation.

\aodprovenance{\mainequation{eq:closure-residual}, \mainequation{eq:shedding-excess}, \mainsection{subsec:support-enclosure-retention-condition}, \mainsection{subsec:yro-odd-window-route-status}.}

\subsection{Complete reference ledger}\label{subsec:complete-elemental-ledger}
The authoritative generated ledger is \path{manual2/data/m1/matter1_elements_1_118_ledger.csv}. It binds each static declaration to its input record, distinguishes labelled elements from unfilled slots, and carries pending admission and unexecuted native status. The support declaration is syntax; the support sum and assigned fields retain the scopes above. Assigned route codes are \texttt{neutral}, \texttt{deficit}, \texttt{surplus} and \texttt{unfilled}; group-boundary and named-anchor flags are independent attributes. The generated table and figure consume the same validated data. Declared compatibility copies must compare equal.

{\scriptsize
\setlength{\tabcolsep}{3.0pt}
\begin{longtable}{@{}rllccrrrrrrll@{}}
\caption{\textbf{Complete declared reference and conditional support ledger.} $Z=1\dots118$ names supplied elements; slots $119\dots126$ are unfilled reference positions. Pressure and route classes are assigned inputs; the support sum is a conditional expression, not an admitted native duration. Grouping follows \mainsection{subsec:coord-zero-support-enclosure-grouping}; the expression specializes \mainequation{eq:curling-support-model}.}\label{tab:matter1-full-ledger}\\
\toprule
$Z$ & Sym & Ref. $A$ & $\mathsf{C}_{2m}$ & Local $a{:}b{:}c$ & $u_-$ & $h$ & $u_+$ & $U$ & $P^{\mathrm{fix}}$ & $C^{\mathrm{fix}}$ & $S^{\mathrm{fix}},X^{\mathrm{fix}}$ & Assigned class\\
\midrule
\endfirsthead
\toprule
$Z$ & Sym & Ref. $A$ & $\mathsf{C}_{2m}$ & Local $a{:}b{:}c$ & $u_-$ & $h$ & $u_+$ & $U$ & $P^{\mathrm{fix}}$ & $C^{\mathrm{fix}}$ & $S^{\mathrm{fix}},X^{\mathrm{fix}}$ & Assigned class\\
\midrule
\endhead
\bottomrule
\endfoot
1 & H & 1 & $\mathsf{C}_{6}$ & $1:0:1$ & 30 & 1 & 4 & 35 & 15 & 18 & $-3,0$ & Deficit\\
2 & He & 4 & $\mathsf{C}_{6}$ & $1:1:0$ & 196 & 1 & 8 & 205 & 18 & 18 & $0,0$ & Neutral\\
3 & Li & 7 & $\mathsf{C}_{6}$ & $2:1:1$ & 362 & 1 & 12 & 375 & 18 & 18 & $0,0$ & Neutral\\
4 & Be & 9 & $\mathsf{C}_{6}$ & $2:2:0$ & 460 & 1 & 16 & 477 & 18 & 18 & $0,0$ & Neutral\\
5 & B & 11 & $\mathsf{C}_{6}$ & $3:2:1$ & 558 & 1 & 20 & 579 & 18 & 18 & $0,0$ & Neutral\\
6 & C & 12 & $\mathsf{C}_{6}$ & $3:3:0$ & 588 & 1 & 24 & 613 & 18 & 18 & $0,0$ & Neutral + anchor\\
7 & N & 14 & $\mathsf{C}_{8}$ & $1:0:1$ & 686 & 1 & 4 & 691 & 24 & 24 & $0,0$ & Neutral\\
8 & O & 16 & $\mathsf{C}_{8}$ & $1:1:0$ & 784 & 1 & 8 & 793 & 21 & 24 & $-3,0$ & Deficit\\
9 & F & 19 & $\mathsf{C}_{8}$ & $2:1:1$ & 950 & 1 & 12 & 963 & 21 & 24 & $-3,0$ & Deficit\\
10 & Ne & 20 & $\mathsf{C}_{8}$ & $2:2:0$ & 980 & 1 & 16 & 997 & 24 & 24 & $0,0$ & Neutral\\
11 & Na & 23 & $\mathsf{C}_{8}$ & $3:2:1$ & 1146 & 1 & 20 & 1167 & 24 & 24 & $0,0$ & Neutral\\
12 & Mg & 24 & $\mathsf{C}_{8}$ & $3:3:0$ & 1176 & 1 & 24 & 1201 & 24 & 24 & $0,0$ & Neutral\\
13 & Al & 27 & $\mathsf{C}_{8}$ & $4:3:1$ & 1342 & 1 & 28 & 1371 & 24 & 24 & $0,0$ & Neutral\\
14 & Si & 28 & $\mathsf{C}_{8}$ & $4:4:0$ & 1372 & 1 & 32 & 1405 & 24 & 24 & $0,0$ & Neutral + anchor\\
15 & P & 31 & $\mathsf{C}_{10}$ & $1:0:1$ & 1538 & 1 & 4 & 1543 & 30 & 30 & $0,0$ & Neutral\\
16 & S & 32 & $\mathsf{C}_{10}$ & $1:1:0$ & 1568 & 1 & 8 & 1577 & 27 & 30 & $-3,0$ & Deficit\\
17 & Cl & 35 & $\mathsf{C}_{10}$ & $2:1:1$ & 1734 & 1 & 12 & 1747 & 27 & 30 & $-3,0$ & Deficit\\
18 & Ar & 40 & $\mathsf{C}_{10}$ & $2:2:0$ & 2036 & 1 & 16 & 2053 & 30 & 30 & $0,0$ & Neutral\\
19 & K & 39 & $\mathsf{C}_{10}$ & $3:2:1$ & 1930 & 1 & 20 & 1951 & 30 & 30 & $0,0$ & Neutral\\
20 & Ca & 40 & $\mathsf{C}_{10}$ & $3:3:0$ & 1960 & 1 & 24 & 1985 & 30 & 30 & $0,0$ & Neutral\\
21 & Sc & 45 & $\mathsf{C}_{10}$ & $4:3:1$ & 2262 & 1 & 28 & 2291 & 30 & 30 & $0,0$ & Neutral\\
22 & Ti & 48 & $\mathsf{C}_{10}$ & $4:4:0$ & 2428 & 1 & 32 & 2461 & 30 & 30 & $0,0$ & Neutral\\
23 & V & 51 & $\mathsf{C}_{10}$ & $5:4:1$ & 2594 & 1 & 36 & 2631 & 30 & 30 & $0,0$ & Neutral\\
24 & Cr & 52 & $\mathsf{C}_{10}$ & $5:5:0$ & 2624 & 1 & 40 & 2665 & 30 & 30 & $0,0$ & Neutral + anchor\\
25 & Mn & 55 & $\mathsf{C}_{12}$ & $1:0:1$ & 2790 & 1 & 4 & 2795 & 36 & 36 & $0,0$ & Neutral\\
26 & Fe & 56 & $\mathsf{C}_{12}$ & $1:1:0$ & 2820 & 1 & 8 & 2829 & 36 & 36 & $0,0$ & Neutral\\
27 & Co & 59 & $\mathsf{C}_{12}$ & $2:1:1$ & 2986 & 1 & 12 & 2999 & 36 & 36 & $0,0$ & Neutral\\
28 & Ni & 58 & $\mathsf{C}_{12}$ & $2:2:0$ & 2880 & 1 & 16 & 2897 & 36 & 36 & $0,0$ & Neutral\\
29 & Cu & 63 & $\mathsf{C}_{12}$ & $3:2:1$ & 3182 & 1 & 20 & 3203 & 36 & 36 & $0,0$ & Neutral\\
30 & Zn & 64 & $\mathsf{C}_{12}$ & $3:3:0$ & 3212 & 1 & 24 & 3237 & 36 & 36 & $0,0$ & Neutral\\
31 & Ga & 69 & $\mathsf{C}_{12}$ & $4:3:1$ & 3514 & 1 & 28 & 3543 & 36 & 36 & $0,0$ & Neutral\\
32 & Ge & 74 & $\mathsf{C}_{12}$ & $4:4:0$ & 3816 & 1 & 32 & 3849 & 36 & 36 & $0,0$ & Neutral\\
33 & As & 75 & $\mathsf{C}_{12}$ & $5:4:1$ & 3846 & 1 & 36 & 3883 & 36 & 36 & $0,0$ & Neutral\\
34 & Se & 80 & $\mathsf{C}_{12}$ & $5:5:0$ & 4148 & 1 & 40 & 4189 & 33 & 36 & $-3,0$ & Deficit\\
35 & Br & 79 & $\mathsf{C}_{12}$ & $6:5:1$ & 4042 & 1 & 44 & 4087 & 33 & 36 & $-3,0$ & Deficit\\
36 & Kr & 84 & $\mathsf{C}_{12}$ & $6:6:0$ & 4344 & 1 & 48 & 4393 & 36 & 36 & $0,0$ & Neutral + anchor\\
37 & Rb & 85 & $\mathsf{C}_{14}$ & $1:0:1$ & 4374 & 1 & 4 & 4379 & 42 & 42 & $0,0$ & Neutral\\
38 & Sr & 88 & $\mathsf{C}_{14}$ & $1:1:0$ & 4540 & 1 & 8 & 4549 & 42 & 42 & $0,0$ & Neutral\\
39 & Y & 89 & $\mathsf{C}_{14}$ & $2:1:1$ & 4570 & 1 & 12 & 4583 & 42 & 42 & $0,0$ & Neutral\\
40 & Zr & 90 & $\mathsf{C}_{14}$ & $2:2:0$ & 4600 & 1 & 16 & 4617 & 42 & 42 & $0,0$ & Neutral\\
41 & Nb & 93 & $\mathsf{C}_{14}$ & $3:2:1$ & 4766 & 1 & 20 & 4787 & 42 & 42 & $0,0$ & Neutral\\
42 & Mo & 98 & $\mathsf{C}_{14}$ & $3:3:0$ & 5068 & 1 & 24 & 5093 & 42 & 42 & $0,0$ & Neutral\\
43 & Tc & 98 & $\mathsf{C}_{14}$ & $4:3:1$ & 5030 & 1 & 28 & 5059 & 45 & 42 & $3,3$ & Surplus\\
44 & Ru & 102 & $\mathsf{C}_{14}$ & $4:4:0$ & 5264 & 1 & 32 & 5297 & 42 & 42 & $0,0$ & Neutral\\
45 & Rh & 103 & $\mathsf{C}_{14}$ & $5:4:1$ & 5294 & 1 & 36 & 5331 & 42 & 42 & $0,0$ & Neutral\\
46 & Pd & 106 & $\mathsf{C}_{14}$ & $5:5:0$ & 5460 & 1 & 40 & 5501 & 42 & 42 & $0,0$ & Neutral\\
47 & Ag & 107 & $\mathsf{C}_{14}$ & $6:5:1$ & 5490 & 1 & 44 & 5535 & 42 & 42 & $0,0$ & Neutral\\
48 & Cd & 114 & $\mathsf{C}_{14}$ & $6:6:0$ & 5928 & 1 & 48 & 5977 & 42 & 42 & $0,0$ & Neutral\\
49 & In & 115 & $\mathsf{C}_{14}$ & $7:6:1$ & 5958 & 1 & 52 & 6011 & 42 & 42 & $0,0$ & Neutral\\
50 & Sn & 120 & $\mathsf{C}_{14}$ & $7:7:0$ & 6260 & 1 & 56 & 6317 & 42 & 42 & $0,0$ & Neutral + anchor\\
51 & Sb & 121 & $\mathsf{C}_{16}$ & $1:0:1$ & 6290 & 1 & 4 & 6295 & 48 & 48 & $0,0$ & Neutral\\
52 & Te & 130 & $\mathsf{C}_{16}$ & $1:1:0$ & 6864 & 1 & 8 & 6873 & 45 & 48 & $-3,0$ & Deficit\\
53 & I & 127 & $\mathsf{C}_{16}$ & $2:1:1$ & 6622 & 1 & 12 & 6635 & 45 & 48 & $-3,0$ & Deficit\\
54 & Xe & 132 & $\mathsf{C}_{16}$ & $2:2:0$ & 6924 & 1 & 16 & 6941 & 48 & 48 & $0,0$ & Neutral\\
55 & Cs & 133 & $\mathsf{C}_{16}$ & $3:2:1$ & 6954 & 1 & 20 & 6975 & 48 & 48 & $0,0$ & Neutral\\
56 & Ba & 138 & $\mathsf{C}_{16}$ & $3:3:0$ & 7256 & 1 & 24 & 7281 & 48 & 48 & $0,0$ & Neutral\\
57 & La & 139 & $\mathsf{C}_{16}$ & $4:3:1$ & 7286 & 1 & 28 & 7315 & 48 & 48 & $0,0$ & Neutral\\
58 & Ce & 140 & $\mathsf{C}_{16}$ & $4:4:0$ & 7316 & 1 & 32 & 7349 & 48 & 48 & $0,0$ & Neutral\\
59 & Pr & 141 & $\mathsf{C}_{16}$ & $5:4:1$ & 7346 & 1 & 36 & 7383 & 48 & 48 & $0,0$ & Neutral\\
60 & Nd & 142 & $\mathsf{C}_{16}$ & $5:5:0$ & 7376 & 1 & 40 & 7417 & 48 & 48 & $0,0$ & Neutral\\
61 & Pm & 145 & $\mathsf{C}_{16}$ & $6:5:1$ & 7542 & 1 & 44 & 7587 & 51 & 48 & $3,3$ & Surplus\\
62 & Sm & 152 & $\mathsf{C}_{16}$ & $6:6:0$ & 7980 & 1 & 48 & 8029 & 48 & 48 & $0,0$ & Neutral\\
63 & Eu & 153 & $\mathsf{C}_{16}$ & $7:6:1$ & 8010 & 1 & 52 & 8063 & 48 & 48 & $0,0$ & Neutral\\
64 & Gd & 158 & $\mathsf{C}_{16}$ & $7:7:0$ & 8312 & 1 & 56 & 8369 & 48 & 48 & $0,0$ & Neutral\\
65 & Tb & 159 & $\mathsf{C}_{16}$ & $8:7:1$ & 8342 & 1 & 60 & 8403 & 48 & 48 & $0,0$ & Neutral\\
66 & Dy & 164 & $\mathsf{C}_{16}$ & $8:8:0$ & 8644 & 1 & 64 & 8709 & 48 & 48 & $0,0$ & Neutral + anchor\\
67 & Ho & 165 & $\mathsf{C}_{18}$ & $1:0:1$ & 8674 & 1 & 4 & 8679 & 54 & 54 & $0,0$ & Neutral\\
68 & Er & 166 & $\mathsf{C}_{18}$ & $1:1:0$ & 8704 & 1 & 8 & 8713 & 54 & 54 & $0,0$ & Neutral\\
69 & Tm & 169 & $\mathsf{C}_{18}$ & $2:1:1$ & 8870 & 1 & 12 & 8883 & 54 & 54 & $0,0$ & Neutral\\
70 & Yb & 174 & $\mathsf{C}_{18}$ & $2:2:0$ & 9172 & 1 & 16 & 9189 & 54 & 54 & $0,0$ & Neutral\\
71 & Lu & 175 & $\mathsf{C}_{18}$ & $3:2:1$ & 9202 & 1 & 20 & 9223 & 54 & 54 & $0,0$ & Neutral\\
72 & Hf & 180 & $\mathsf{C}_{18}$ & $3:3:0$ & 9504 & 1 & 24 & 9529 & 54 & 54 & $0,0$ & Neutral\\
73 & Ta & 181 & $\mathsf{C}_{18}$ & $4:3:1$ & 9534 & 1 & 28 & 9563 & 54 & 54 & $0,0$ & Neutral\\
74 & W & 184 & $\mathsf{C}_{18}$ & $4:4:0$ & 9700 & 1 & 32 & 9733 & 54 & 54 & $0,0$ & Neutral\\
75 & Re & 187 & $\mathsf{C}_{18}$ & $5:4:1$ & 9866 & 1 & 36 & 9903 & 54 & 54 & $0,0$ & Neutral\\
76 & Os & 192 & $\mathsf{C}_{18}$ & $5:5:0$ & 10168 & 1 & 40 & 10209 & 54 & 54 & $0,0$ & Neutral\\
77 & Ir & 193 & $\mathsf{C}_{18}$ & $6:5:1$ & 10198 & 1 & 44 & 10243 & 54 & 54 & $0,0$ & Neutral\\
78 & Pt & 195 & $\mathsf{C}_{18}$ & $6:6:0$ & 10296 & 1 & 48 & 10345 & 54 & 54 & $0,0$ & Neutral\\
79 & Au & 197 & $\mathsf{C}_{18}$ & $7:6:1$ & 10394 & 1 & 52 & 10447 & 54 & 54 & $0,0$ & Neutral\\
80 & Hg & 202 & $\mathsf{C}_{18}$ & $7:7:0$ & 10696 & 1 & 56 & 10753 & 54 & 54 & $0,0$ & Neutral\\
81 & Tl & 205 & $\mathsf{C}_{18}$ & $8:7:1$ & 10862 & 1 & 60 & 10923 & 54 & 54 & $0,0$ & Neutral\\
82 & Pb & 208 & $\mathsf{C}_{18}$ & $8:8:0$ & 11028 & 1 & 64 & 11093 & 54 & 54 & $0,0$ & Neutral\\
83 & Bi & 209 & $\mathsf{C}_{18}$ & $9:8:1$ & 11058 & 1 & 68 & 11127 & 54 & 54 & $0,0$ & Neutral\\
84 & Po & 209 & $\mathsf{C}_{18}$ & $9:9:0$ & 11020 & 1 & 72 & 11093 & 57 & 54 & $3,3$ & Surplus + anchor\\
85 & At & 210 & $\mathsf{C}_{20}$ & $1:0:1$ & 11050 & 1 & 4 & 11055 & 63 & 60 & $3,3$ & Surplus\\
86 & Rn & 222 & $\mathsf{C}_{20}$ & $1:1:0$ & 11828 & 1 & 8 & 11837 & 63 & 60 & $3,3$ & Surplus\\
87 & Fr & 223 & $\mathsf{C}_{20}$ & $2:1:1$ & 11858 & 1 & 12 & 11871 & 63 & 60 & $3,3$ & Surplus\\
88 & Ra & 226 & $\mathsf{C}_{20}$ & $2:2:0$ & 12024 & 1 & 16 & 12041 & 63 & 60 & $3,3$ & Surplus\\
89 & Ac & 227 & $\mathsf{C}_{20}$ & $3:2:1$ & 12054 & 1 & 20 & 12075 & 63 & 60 & $3,3$ & Surplus\\
90 & Th & 232 & $\mathsf{C}_{20}$ & $3:3:0$ & 12356 & 1 & 24 & 12381 & 66 & 60 & $6,6$ & Surplus\\
91 & Pa & 231 & $\mathsf{C}_{20}$ & $4:3:1$ & 12250 & 1 & 28 & 12279 & 66 & 60 & $6,6$ & Surplus\\
92 & U & 238 & $\mathsf{C}_{20}$ & $4:4:0$ & 12688 & 1 & 32 & 12721 & 66 & 60 & $6,6$ & Surplus\\
93 & Np & 237 & $\mathsf{C}_{20}$ & $5:4:1$ & 12582 & 1 & 36 & 12619 & 66 & 60 & $6,6$ & Surplus\\
94 & Pu & 244 & $\mathsf{C}_{20}$ & $5:5:0$ & 13020 & 1 & 40 & 13061 & 66 & 60 & $6,6$ & Surplus\\
95 & Am & 243 & $\mathsf{C}_{20}$ & $6:5:1$ & 12914 & 1 & 44 & 12959 & 66 & 60 & $6,6$ & Surplus\\
96 & Cm & 247 & $\mathsf{C}_{20}$ & $6:6:0$ & 13148 & 1 & 48 & 13197 & 69 & 60 & $9,9$ & Surplus\\
97 & Bk & 247 & $\mathsf{C}_{20}$ & $7:6:1$ & 13110 & 1 & 52 & 13163 & 69 & 60 & $9,9$ & Surplus\\
98 & Cf & 251 & $\mathsf{C}_{20}$ & $7:7:0$ & 13344 & 1 & 56 & 13401 & 69 & 60 & $9,9$ & Surplus\\
99 & Es & 252 & $\mathsf{C}_{20}$ & $8:7:1$ & 13374 & 1 & 60 & 13435 & 69 & 60 & $9,9$ & Surplus\\
100 & Fm & 257 & $\mathsf{C}_{20}$ & $8:8:0$ & 13676 & 1 & 64 & 13741 & 69 & 60 & $9,9$ & Surplus\\
101 & Md & 258 & $\mathsf{C}_{20}$ & $9:8:1$ & 13706 & 1 & 68 & 13775 & 69 & 60 & $9,9$ & Surplus\\
102 & No & 259 & $\mathsf{C}_{20}$ & $9:9:0$ & 13736 & 1 & 72 & 13809 & 72 & 60 & $12,12$ & Surplus\\
103 & Lr & 266 & $\mathsf{C}_{20}$ & $10:9:1$ & 14174 & 1 & 76 & 14251 & 72 & 60 & $12,12$ & Surplus\\
104 & Rf & 267 & $\mathsf{C}_{20}$ & $10:10:0$ & 14204 & 1 & 80 & 14285 & 72 & 60 & $12,12$ & Surplus + anchor\\
105 & Db & 268 & $\mathsf{C}_{22}$ & $1:0:1$ & 14234 & 1 & 4 & 14239 & 78 & 66 & $12,12$ & Surplus\\
106 & Sg & 269 & $\mathsf{C}_{22}$ & $1:1:0$ & 14264 & 1 & 8 & 14273 & 78 & 66 & $12,12$ & Surplus\\
107 & Bh & 270 & $\mathsf{C}_{22}$ & $2:1:1$ & 14294 & 1 & 12 & 14307 & 78 & 66 & $12,12$ & Surplus\\
108 & Hs & 277 & $\mathsf{C}_{22}$ & $2:2:0$ & 14732 & 1 & 16 & 14749 & 81 & 66 & $15,15$ & Surplus\\
109 & Mt & 278 & $\mathsf{C}_{22}$ & $3:2:1$ & 14762 & 1 & 20 & 14783 & 81 & 66 & $15,15$ & Surplus\\
110 & Ds & 281 & $\mathsf{C}_{22}$ & $3:3:0$ & 14928 & 1 & 24 & 14953 & 81 & 66 & $15,15$ & Surplus\\
111 & Rg & 282 & $\mathsf{C}_{22}$ & $4:3:1$ & 14958 & 1 & 28 & 14987 & 81 & 66 & $15,15$ & Surplus\\
112 & Cn & 285 & $\mathsf{C}_{22}$ & $4:4:0$ & 15124 & 1 & 32 & 15157 & 81 & 66 & $15,15$ & Surplus\\
113 & Nh & 286 & $\mathsf{C}_{22}$ & $5:4:1$ & 15154 & 1 & 36 & 15191 & 81 & 66 & $15,15$ & Surplus\\
114 & Fl & 289 & $\mathsf{C}_{22}$ & $5:5:0$ & 15320 & 1 & 40 & 15361 & 84 & 66 & $18,18$ & Surplus\\
115 & Mc & 290 & $\mathsf{C}_{22}$ & $6:5:1$ & 15350 & 1 & 44 & 15395 & 84 & 66 & $18,18$ & Surplus\\
116 & Lv & 293 & $\mathsf{C}_{22}$ & $6:6:0$ & 15516 & 1 & 48 & 15565 & 84 & 66 & $18,18$ & Surplus\\
117 & Ts & 294 & $\mathsf{C}_{22}$ & $7:6:1$ & 15546 & 1 & 52 & 15599 & 84 & 66 & $18,18$ & Surplus\\
118 & Og & 294 & $\mathsf{C}_{22}$ & $7:7:0$ & 15508 & 1 & 56 & 15565 & 84 & 66 & $18,18$ & Surplus\\
119 & -- & 297 & $\mathsf{C}_{22}$ & $8:7:1$ & 15674 & 1 & 60 & 15735 & 69 & 66 & $3,3$ & Unfilled fixture\\
120 & -- & 298 & $\mathsf{C}_{22}$ & $8:8:0$ & 15704 & 1 & 64 & 15769 & 72 & 66 & $6,6$ & Unfilled fixture\\
121 & -- & 301 & $\mathsf{C}_{22}$ & $9:8:1$ & 15870 & 1 & 68 & 15939 & 75 & 66 & $9,9$ & Unfilled fixture\\
122 & -- & 302 & $\mathsf{C}_{22}$ & $9:9:0$ & 15900 & 1 & 72 & 15973 & 78 & 66 & $12,12$ & Unfilled fixture\\
123 & -- & 305 & $\mathsf{C}_{22}$ & $10:9:1$ & 16066 & 1 & 76 & 16143 & 81 & 66 & $15,15$ & Unfilled fixture\\
124 & -- & 306 & $\mathsf{C}_{22}$ & $10:10:0$ & 16096 & 1 & 80 & 16177 & 84 & 66 & $18,18$ & Unfilled fixture\\
125 & -- & 309 & $\mathsf{C}_{22}$ & $11:10:1$ & 16262 & 1 & 84 & 16347 & 87 & 66 & $21,21$ & Unfilled fixture\\
126 & -- & 310 & $\mathsf{C}_{22}$ & $11:11:0$ & 16292 & 1 & 88 & 16381 & 90 & 66 & $24,24$ & Unfilled fixture\\
\end{longtable}
}

